\chapter{Concepts}\label{ch:concepts}

This chapter explains the ideas that the whole interface relies on: how points are passed, what
a NURBS curve and surface are, how parameters, knots and weights behave, which conventions the
constructors follow, how tolerances are handled, and how errors, immutability, threads and units
work. Later chapters refer back to these sections.

\section{Points and vectors}\label{sec:concepts-vec3}

Points and vectors are represented by the class \apiref{Vec3} with the attributes \code{x},
\code{y} and \code{z}. It supports the usual vector arithmetic (\code{+}, \code{-},
multiplication by a scalar from either side), iteration (\code{list(v)}, unpacking
\code{x, y, z = v}) and conversion to a tuple with \code{to_tuple()}. The free functions
\apiref{dot}, \apiref{cross}, \apiref{norm} and \apiref{distance} complete the set.

\paragraph{Implicit conversion.} Wherever a function expects a point or a vector, any sequence
of exactly three numbers is accepted: a tuple, a list, a \code{Vec3} or a one-dimensional
NumPy array, including a row of an $(n,3)$ array. Results are always returned as \code{Vec3}.
A sequence of the wrong length is rejected with a \code{TypeError} that lists the accepted
signatures. There are no two-dimensional points: points in a plane are given with $z=0$.

\paragraph{Bulk data.} Functions that take or return many points use NumPy arrays: point lists
of shape $(n,3)$ (fitting, polylines), point grids of shape $(n_u,n_v,3)$ (surface fitting) and
the results of \code{curve.sample(n)} and \code{surface.sample(nu, nv)}. Control points are
returned as arrays too (\code{curve.control_points}, shape $(n,3)$).

\example{04_points_and_vectors.py}

\begin{lstlisting}[style=text]
a + b = Vec3(1.000000, 3.000000, 2.000000)  a - b = Vec3(1.000000, 1.000000, 2.000000)  2 * a = Vec3(2.000000, 4.000000, 4.000000)
components: 1.0 2.0 2.0  tuple: (1.0, 2.0, 2.0)  list: [1.0, 2.0, 2.0]
dot 2.0  cross Vec3(-2.000000, 0.000000, 1.000000)  norm 3.0  distance 2.449489742783178
length of line2: 5.000000000000001
midpoint Vec3(1.500000, 2.000000, 0.000000) -> [1.5 2.  0. ]
sample(5): (5, 3) [0.   0.75 1.5  2.25 3.  ]
\end{lstlisting}

\begin{note}[Printed precision]
The text representation of \code{Vec3} (and of \code{Tolerance}) shows six decimal places. The
stored values are full double-precision numbers. Print the components with a format string,
for example \code{f"{p.x:.15g}"}, when you need all digits.
\end{note}

\section{NURBS curves}\label{sec:concepts-nurbs}

\subsection{Definition}
A NURBS curve of degree $p$ with control points $\mathbf P_0,\dots,\mathbf P_n$, positive weights
$w_0,\dots,w_n$ and a knot vector $U=\{u_0,\dots,u_{n+p+1}\}$ is
\begin{equation}
  \mathbf C(u) = \frac{\sum_{i=0}^{n} N_{i,p}(u)\, w_i\, \mathbf P_i}
                      {\sum_{i=0}^{n} N_{i,p}(u)\, w_i},
  \qquad u \in [u_p, u_{n+1}] .
\end{equation}
The $N_{i,p}$ are the B-spline basis functions, defined recursively by the Cox--de Boor formula:
\begin{align}
  N_{i,0}(u) &= \begin{cases} 1 & u_i \le u < u_{i+1} \\ 0 & \text{otherwise,} \end{cases}\\
  N_{i,p}(u) &= \frac{u-u_i}{u_{i+p}-u_i}\,N_{i,p-1}(u)
              + \frac{u_{i+p+1}-u}{u_{i+p+1}-u_{i+1}}\,N_{i+1,p-1}(u)
\end{align}
(with $0/0 := 0$). Each $N_{i,p}$ is a piecewise polynomial of degree $p$. It is nonzero only on
$[u_i, u_{i+p+1})$, the basis functions are nonnegative and they sum to one. Hence:
\begin{itemize}
  \item \textbf{Local control.} Moving $\mathbf P_i$ changes the curve only for
    $u\in[u_i,u_{i+p+1})$.
  \item \textbf{Convex hull.} On each knot span the curve lies in the convex hull of the $p+1$
    active control points (for positive weights).
  \item \textbf{Affine invariance.} Transforming the control points transforms the curve. This is
    how \code{transformed()} works.
\end{itemize}
Internally, evaluation uses homogeneous coordinates $\mathbf P^w_i = (w_i\mathbf P_i,\, w_i)$:
the curve is a polynomial B-spline in four dimensions, projected back by division by the last
coordinate. \NF{} implements the algorithms of Piegl and Tiller, \emph{The NURBS Book}, and names
them in the source code (for example A2.1 \emph{FindSpan}, A4.2 \emph{RatCurveDerivs}, A5.1
\emph{CurveKnotIns}, A7.1 \emph{MakeNurbsCircle}).

\subsection{Knot vectors}
A knot vector is non-decreasing and has exactly $n_{\text{ctrl}}+p+1$ entries. The number of
times a value appears is its \emph{multiplicity}. At an interior knot of multiplicity $m$ the
curve is $C^{p-m}$ continuous. \NF{} requires $m\le p$ for interior knots, so curves are at
least $C^0$ continuous.

A knot vector is \emph{clamped} when its first and last values have multiplicity $p+1$. Then the
curve starts at the first control point and ends at the last one, tangent to the control
polygon. All constructors create clamped curves. Unclamped curves are allowed, and
\apiref{NurbsCurve.clamped} converts them.

Helpers for knot vectors:
\begin{itemize}
  \item \apiref{clamped_uniform_knots} \code{(degree, n_ctrl)}: a clamped vector with uniform
    interior knots, e.g.\ \code{clamped_uniform_knots(3, 6)} $=[0,0,0,0,\frac13,\frac23,1,1,1,1]$.
  \item \apiref{unique_knots} \code{(knots)}: the distinct values and their multiplicities,
    e.g.\ \code{([0.0, 0.5, 1.0], [3, 2, 3])}.
  \item \apiref{validate_knots} \code{(knots, degree, n_ctrl)}: raises
    \code{GeometryError} if the vector is invalid.
\end{itemize}

\subsection{Weights and rational curves}
With all weights equal, the denominator is constant and the curve is a polynomial B-spline.
Different weights make it \emph{rational}, and only rational curves can represent conics
exactly. Increasing $w_i$ pulls the curve towards $\mathbf P_i$. Multiplying all weights by
the same factor does not change the curve. \code{curve.is_rational} reports whether the weights
differ.

\subsection{Parameter domain}
The curve is defined on $[u_p, u_{n+1}]$, returned as \code{curve.domain}. All constructors use
$[0,1]$. \apiref{reparametrize} maps a curve to another interval without changing its shape.
The parameter is generally \emph{not} proportional to arc length. Use \apiref{arc_length}
and \apiref{parameter_at_length} to convert between parameter and length.

\subsection{Building a curve from its definition}
Usually curves come from constructors, but \apiref{NurbsCurve} can also be created from its
data. The following script builds a quarter circle as a rational quadratic with the middle
weight $\cos 45^\circ$, shows that the same control points without weights give a parabola,
inserts a knot and computes derivatives:

\example{04_nurbs_by_hand.py}

\begin{lstlisting}[style=text]
<NurbsCurve degree=2 n_ctrl=3 rational> rational: True clamped: True
max | |C(u)| - 1 | = 2.22e-16
non-rational: |C(0.5)| = 1.060660
knots after inserting 0.3 twice: [0.0, 0.0, 0.0, 0.3, 0.3, 1.0, 1.0, 1.0]
max deviation after insertion: 2.5e-16
C(0) = Vec3(1.000000, 0.000000, 0.000000)  C'(0) = Vec3(0.000000, 1.414214, 0.000000)  C''(0) = Vec3(-2.000000, 0.828427, 0.000000)
curvature at u = 0.5: 1.000000000000
GeometryError: knot vector size must be n_ctrl + degree + 1 = 6 (got 4)
\end{lstlisting}

The constructor signature is
\code{NurbsCurve(degree, knots, control_points, weights=[])}. An empty weight list means a
non-rational curve. \code{NurbsCurve.from_points(degree, points)} creates a clamped uniform
B-spline that uses the given points as its control points (it does \emph{not} interpolate them;
for interpolation see \cref{ch:fitting}).

\subsection{Operations that change the representation only}
Several operations change knots, degree or the number of control points but leave the geometric
curve unchanged (up to rounding): knot insertion and refinement
(\apiref{NurbsCurve.insert_knot}, \apiref{NurbsCurve.refine_knots}), degree elevation
(\apiref{NurbsCurve.elevate_degree}), Bézier decomposition
(\apiref{NurbsCurve.decompose_bezier}) and clamping. Knot removal and degree reduction change the
curve within a given tolerance, measured as a Cartesian deviation in model units. These
operations are used to make curves \emph{compatible} (same degree and knot vector), which loft,
ruled and Gordon surfaces need. \Cref{ch:curves} describes them in detail.

\section{NURBS surfaces}\label{sec:concepts-surfaces}

A tensor-product NURBS surface of degrees $(p,q)$ has a net of $(n+1)\times(m+1)$ control points
$\mathbf P_{ij}$ with weights $w_{ij}$ and two knot vectors $U$ and $V$:
\begin{equation}
  \mathbf S(u,v) = \frac{\sum_{i}\sum_{j} N_{i,p}(u)\,N_{j,q}(v)\,w_{ij}\,\mathbf P_{ij}}
                        {\sum_{i}\sum_{j} N_{i,p}(u)\,N_{j,q}(v)\,w_{ij}} .
\end{equation}
For fixed $v$ it is a curve in $u$ (an \emph{iso-curve}, \apiref{NurbsSurface.iso_curve_v}), and
vice versa.

\subsection{Control net ordering}
The control points are stored \textbf{row-major with $i$ along $u$}: point $\mathbf P_{ij}$ has
the index $i\cdot n_v + j$, where $n_v$ is \code{num_v}. The constructor is
\code{NurbsSurface(degree_u, degree_v, knots_u, knots_v, num_u, num_v, control_points,
weights=[])} with \code{control_points} as a flat list or an array of shape
$(n_u n_v, 3)$. \code{surface.control_points} returns the same layout.

\example{04_surface_net.py}

\begin{lstlisting}[style=text]
S(1, 0) = Vec3(2.000000, 0.000000, 0.000000)  S(0, 1) = Vec3(0.000000, 1.000000, 0.000000)  S(1, 1) = Vec3(2.000000, 1.000000, 1.000000)
net shape 2 x 2  control_points: (4, 3)
normal at (0.5, 0.5): Vec3(-0.218218, -0.436436, 0.872872)
S_u = Vec3(2.000000, 0.000000, 0.500000)  S_v = Vec3(0.000000, 1.000000, 0.500000)
transposed normal: Vec3(0.218218, 0.436436, -0.872872)
\end{lstlisting}

\begin{warning}[gmsh and other libraries]
Other libraries may order the net differently. gmsh, for example, expects $u$ to run fastest.
The \code{nurbforge.mesh} module reorders automatically. Reorder yourself only when you pass raw
control points to other software.
\end{warning}

\subsection{Normals and orientation}
The normal of a surface is $\mathbf S_u\times\mathbf S_v$, normalised
(\apiref{NurbsSurface.normal}). At degenerate points such as the poles of a revolved sphere,
where $\mathbf S_u$ vanishes, \NF{} returns the limit normal. Swapping $u$ and $v$
(\apiref{NurbsSurface.transposed}) or reversing one direction flips the normal. The orientation
matters for offsets, fillets, trimmed faces and shells: a closed shell is \emph{oriented} when
all face normals point out of the solid.

\section{Conventions}\label{sec:concepts-conventions}

\begin{longtable}{@{}p{0.27\textwidth}p{0.67\textwidth}@{}}
\toprule
\textbf{Item} & \textbf{Convention} \\
\midrule
\endhead
Knot vectors & Non-decreasing, size \code{n_ctrl + degree + 1}, interior multiplicity
  $\le$ degree. Clamped vectors are typical but not required. \\
Default domain & $[0,1]$ for all constructors (primitives, construction, fitting). \\
Surface net & Row-major, \code{index = i * num_v + j}, $i$ along $u$. \\
Revolution & $u$ = circumferential (degree 2, rational), $v$ = profile. \\
Extrusion, ruled & $u$ = profile, $v$ = sweep direction (degree 1). \\
Loft & $u$ along the sections, $v$ across them. \\
Angles & Radians in the C++ and Python API; \textbf{degrees} in the LLM tools. \\
Arcs & \code{make_arc(center, x_axis, y_axis, radius, start, end)} runs from angle
  \code{start} to \code{end}, measured from \code{x_axis} towards \code{y_axis}. \\
Circles & Counter-clockwise about the given normal. The start point lies in a direction
  perpendicular to the normal chosen by the kernel; for the normal $+z$ it is on the $+y$ axis,
  so $u=0.25$ is on $-x$. Use \code{make_circle_3pt} or \code{make_arc} to control the start. \\
Trim loops & In the $(u,v)$ parameter space of the surface; outer loop counter-clockwise, holes
  clockwise. \\
Shells & Normals point outwards (away from the material); voids have normals pointing into the
  void. \\
Planar sections & Loops oriented with the solid on the left: outer loops counter-clockwise,
  holes clockwise in the plane coordinates. \\
Tolerances & \code{Tolerance(length=1e-10, parametric=1e-10, angular=1e-9)}; see below. \\
\bottomrule
\end{longtable}

\section{Exact and approximated constructions}\label{sec:concepts-exact}

Every operation in \NF{} falls into one of three classes:

\begin{description}[style=nextline,leftmargin=1.5em]
  \item[Exact] The result represents the mathematical object exactly; the only error is
    floating-point rounding (relative $10^{-15}$ to $10^{-13}$). Lines, polylines, polygons,
    arcs and circles, conics, extrusions, revolutions, ruled surfaces, bilinear patches, the
    solid primitives, translational sweeps, transforms,
    knot insertion and refinement, degree elevation, splitting, and fillet arcs between curves.
  \item[Approximated to a verified tolerance] The exact result is not a NURBS (or not one of
    reasonable size). The operation builds an approximation and checks it at dense check
    points, refining until the requested tolerance is met: helices and spirals, general
    sweeps, offsets, surface fillets, aspheric and freeform optical surfaces, sag surfaces,
    Coons and Gordon surfaces (where the input curves must be made compatible),
    intersection and section curves on general surfaces, trim curves mapped to 3D, fits to a
    tolerance, degree reduction and knot removal.
  \item[Numerical results] Scalars computed to a relative accuracy: lengths, areas, volumes,
    mass properties (adaptive Gauss--Kronrod integration with \code{rel_tol}), closest points
    and intersection points (Newton iterations to the \code{Tolerance}).
\end{description}

\section{The tolerance model}\label{sec:concepts-tolerances}

\subsection{Three kinds of tolerance}
Lengths, parameter values and angles have different units and scale differently. \NF{}
therefore separates them (architecture decision ADR~0005):

\begin{tabularx}{\textwidth}{@{}llX@{}}
\toprule
\textbf{Field} & \textbf{Default} & \textbf{Used for} \\
\midrule
\code{length} & $10^{-10}$ (model units) & Zero lengths, coincident points, parallel lines,
  convergence of point iterations. \\
\code{parametric} & $10^{-10}$ (parameter units) & Equal knots, multiplicities, snapping to
  parameter values. \\
\code{angular} & $10^{-9}$ (radians) & Angle limits, orthogonality and parallelism of unit
  vectors. \\
\bottomrule
\end{tabularx}

The defaults are the kernel constants \code{nf::tol::kLength}, \code{kParametric} and
\code{kAngular} (C++). Weights are compared \emph{relatively}, because a NURBS does not change
when all weights are scaled.

Iterative algorithms take an optional \apiref{Tolerance} object:
\code{nf.closest_point(curve, p, nf.Tolerance(length=1e-6))}. Construct it with keyword arguments.
Fields not given keep their defaults.

\subsection{Tolerances of approximating operations}
Functions that build approximations take a scalar \code{tolerance} in model units, for example
\code{offset_surface(s, d, tolerance=1e-7)} or \code{SweepOptions(tolerance=1e-6)}. The meaning
is always the same: the maximum deviation of the result from the exact object, \emph{checked} at
dense check points (typically eight per knot span, or $8\times8$ per patch), with knots inserted
until the check passes. The tolerance is verified at these points, not proven in between. The
test suite evaluates the results at additional random points and finds errors within 1.5 to 2
times the tolerance. For a strict bound, request a smaller tolerance.

\example{04_tolerances.py}

\begin{lstlisting}[style=text]
defaults: length 1e-10, parametric 1e-10, angular 1e-09
repr rounds to 6 digits: Tolerance(length=0.000000, parametric=0.000000, angular=0.000000)
closest point u = 0.750000, distance 2.000000000000
helix: degree 5, 216 control points, max deviation 7.6e-09 (tolerance 1e-8)
\end{lstlisting}

\begin{warning}[The repr of Tolerance]
\code{print(nf.Tolerance())} shows zeros because the representation rounds to six decimal
places. The actual values are $10^{-10}$ and $10^{-9}$, as the attributes show.
\end{warning}

\subsection{Choosing tolerances}
\begin{itemize}
  \item Choose approximation tolerances relative to the size of the model: for a part of size
    100\,mm, $10^{-6}$ (one nanometre) is far below any manufacturing tolerance. Tighter
    tolerances make the results larger (more knots and control points) and slower to compute.
  \item Many options use \code{0} as ``automatic''. For example, a \code{fit_tolerance} of 0
    in \code{SurfaceIntersectionOptions} means $10^{-7}$ of the model size.
    \Cref{app:tolerances} lists every such parameter with its default.
  \item The sewing tolerance (\code{Shell.sew(faces, tolerance=1e-6)}) and the Boolean tolerance
    (\code{boolean(a, b, op, tolerance=1e-6)}) are distances in model units. For models in
    metres or micrometres scale them accordingly (\cref{sec:concepts-units}).
\end{itemize}

\section{Errors}\label{sec:concepts-errors}

\begin{description}[style=nextline,leftmargin=1.5em]
  \item[\code{nf.GeometryError}] A subclass of \code{ValueError}, raised for invalid input or
    when a geometric operation is impossible: a knot vector of the wrong size, a zero axis, an
    offset that reaches a radius of curvature, a Boolean of an open shell, a tolerance that
    cannot be met within the allowed refinement. The message names the function and the
    violated condition, often with the offending values.
  \item[\code{TypeError}] Raised by the bindings when arguments have the wrong type, for example
    a point with two coordinates. The message lists the accepted signatures.
  \item[\code{NotImplementedError}] Reserved for recognised requests whose implementation is
    still planned.
  \item[\code{nurbforge.mesh.MeshTopologyError}] A subclass of \code{ValueError} for surface
    meshes that are not watertight or orientable (\cref{ch:meshing}).
\end{description}

Since \code{GeometryError} is a \code{ValueError}, generic handlers keep working. Catch
\code{nf.GeometryError} to distinguish geometric failures from other value errors. Errors never
leave partially modified objects behind: inputs are immutable, and a failed operation returns
nothing.

\section{Immutability and value semantics}\label{sec:concepts-immutability}

Curves, surfaces, trimmed surfaces and shells are immutable. Methods such as
\code{insert_knot}, \code{transformed}, \code{reversed}, \code{oriented} or
\code{with_face_labels} return \emph{new} objects:

\begin{lstlisting}
c2 = c.insert_knot(0.5)        # c is unchanged
moved = shell.transformed(nf.Transform.translation((0, 0, 10)))
\end{lstlisting}

Calling \code{c.insert_knot(0.5)} without using the return value therefore has no effect. In
return, objects can be stored in containers, shared between parts of a model and reused without
defensive copies. The geometry objects cannot be pickled. To move them between processes, write
them to STEP text with \apiref{to_step} and read them back with \apiref{from_step}, or rebuild
them from the script that created them.

\section{Thread safety}\label{sec:concepts-threads}

All \code{const} queries may run concurrently on shared objects: evaluation, derivatives, closest
points, intersections, sections, critical heights, lengths, areas, volumes and mass properties,
and STEP reading and writing. Internal caches, such as trim polylines, are built in the
constructors, so reading never writes. The C++ test suite runs these queries concurrently under
ThreadSanitizer.

In Python, most functions hold the global interpreter lock while they compute. Threads therefore
give no speed-up, except for \apiref{hatch} and STEP reading, which release the lock. The
surface/surface intersection uses several threads internally. For parallel batch work in Python
use processes and pass geometry as STEP text or recreate it in each process.

\section{Units}\label{sec:concepts-units}

The kernel works in dimensionless \emph{model units}. It is up to you whether one unit is a
millimetre, a micrometre or a metre. Units enter in three places only:

\begin{enumerate}
  \item \textbf{Tolerances} given in model units (\code{length}, sewing and Boolean tolerances,
    fit tolerances) must suit the scale. The defaults ($10^{-10}$ for comparisons, $10^{-6}$ for
    sewing and Booleans, $10^{-7}$ for most fits) suit models of size 1 to 1000.
  \item \textbf{STEP files} declare a unit: \code{write_step(..., unit="mm")} writes the
    coordinates unchanged and states what one model unit means. \code{read_step(...,
    unit="mm")} converts the file's coordinates to the requested unit (\cref{ch:step}).
  \item \textbf{gmsh} uses the coordinates as they are. Mesh sizes are in model units.
\end{enumerate}

For a microlens in micrometres, model it in micrometres, use tolerances such as $10^{-6}$\,µm
or larger, and write the STEP file with \code{unit="um"}.
